\section{F検定/分散分析の基礎}
\label{b19_Ftest}

\subsection{授業内容}
\subsubsection{概要}
前回学んだ群間計画(Between-subjects Design)の理論的基礎として，平方和の分解について理解した。これと統計的判断すなわちF検定を組み合わせることで，実際のデータに基づいて効果の有意性を評価できるようになる。Rを用いて一要因および二要因の群間計画データを実際に分析し，分散分析表の読み方や結果の解釈方法を習得する。基本的な分析には\verb|aov()|関数と\verb|lm()|関数を用い，それらの出力結果からSS(平方和)，df(自由度)，MS(平均平方)，F値などの情報をどのように読み取るかを学ぶ。また，検定の帰無仮説と対立仮説を改めて確認し，有意な結果が得られた場合に，どの群間に具体的な差があるのかを特定するための下位検定(post-hoc test)の方法についても学ぶ。

\subsubsection{コマ主題細目}
\begin{description}

	\item[F検定とその解釈] 平方和を分解した後，効果の統計的有意性を評価するためにF検定を用いる。F値は，効果の平均平方($MS = SS/df$)を誤差の平均平方で割った値(F = MS効果/MS誤差)として算出される。自由度(df)の概念と計算方法，F分布の特性とF値の解釈について学ぶ。特に，F値が大きいほど効果が誤差に比べて相対的に大きいことを意味し，これが統計的に有意であるかどうかを判断するための基準となることを理解する。また，多要因計画における主効果と交互作用それぞれのF検定と，その結果の解釈方法についても学ぶ。
	\begin{flushright}
		$\to$ \textcite{Kawabata201408}のP.125--130，\textcite{Yamada2004}のP.190--196.
	\end{flushright}

	\item[線形モデルとしてみた分散分析] 分散分析を線形モデルの一種として理解する。平方和の分解のモデルが回帰分析と同じ形をしていること，帰無仮説を数式で表現すると水平な直線であり，対立仮説は傾きのある直線として表現できることを確認する。これにより，分散分析が回帰分析の特別なケースであることを理解し，Rの\verb|lm()|関数を用いて分散分析を実行できるようになる。

	\item[データの準備と基本操作] 分析に先立って，適切なデータの読み込み方法と前処理を学ぶ。特にRでの分散分析では，カテゴリカル変数を因子型(factor)として扱うことが重要である。csvファイルやExcelファイルからのデータの読み込み方法，データフレームの構造確認，因子型への変換方法，データの要約統計量の算出方法などの基本操作を復習する。また，欠損値の処理方法についても学び，分析前の適切なデータチェックの重要性を理解する。
	      \begin{flushright}
		      $\to$ \textcite{kosugi2019}のP.46--70，\textcite{Kinosady2021}のP.28--42.
	      \end{flushright}

	\item[一要因群間計画の分析] Rの\verb|aov()|関数と\verb|lm()|関数を用いた一要因分散分析の実行方法を学ぶ。公式表記(formula notation)の書き方(例：\verb|y ~ group|)，関数の使い方，出力結果の読み方について理解する。\verb|summary()|関数を用いた分散分析表(ANOVA table)の表示方法と，そこに含まれるSS(平方和)，df(自由度)，MS(平均平方)，F値，p値の意味を確認する。
	      \begin{flushright}
		      $\to$ \textcite{kosugi2019}のP.146--152，\textcite{Hashimoto201607}のP.26--55.
	      \end{flushright}

	\item[下位検定]
	分散分析における下位検定(post-hoc test)の必要性と目的について学ぶ。分散分析における帰無仮説と対立仮説を確認し，何がどのように判断されたのかを改めて確認する。有意な結果が得られた場合，どの群間に具体的な差があるのかを特定するために下位検定を行う必要があることを理解する。代表的な下位検定の方法(例：TukeyのHSD，Bonferroni法，Scheffé法など)について理解する。また，\verb|TukeyHSD()|関数を用いた多重比較の方法や，効果量($\eta^2$，$\omega^2$など)の算出方法についても学ぶ。実際のデータセットを用いた演習を通じて，一要因群間計画の分析手順と結果の解釈を身につける。



\end{description}
\subsubsection{キーワード}
\begin{itemize}
	\item \verb|aov()|関数と\verb|lm()|関数
	\item 分散分析表の読み方(SS, df, MS, F値)
	\item 多重比較と下位検定
\end{itemize}
\subsection{授業情報}
\paragraph{コマの展開方法}
演習
\paragraph{標準シラバスにおける位置づけ}
科目番号5 心理学統計法；2統計に関する基礎的な知識；C/D/E エクセル，R，SPSS入門
\subsubsection{予習・復習課題}
\paragraph{予習}
前回学んだ群間計画(Between Design)の基本概念と分散分析の理論を復習しておくこと。特に，平方和の分解，F値の意味，主効果と交互作用の概念について理解を確認しておくとよい。また，RとRStudioの基本操作，特にデータの読み込み方法やデータフレームの操作方法について再確認しておくこと。必要に応じて，第3回，第4回の授業で学んだR操作についての資料を見直しておくとよい。

\paragraph{復習}
授業で学んだ内容を定着させるために，以下の課題に取り組むこと：

\begin{enumerate}
\item 提供されたデータセットを用いて，一要因群間計画の分散分析を実行し，結果を適切に解釈する
\item 提供されたデータセットを用いて，二要因群間計画の分散分析を実行し，主効果と交互作用の結果を解釈する
\item \verb|TukeyHSD()|関数を用いて多重比較を行い，どの群間に有意差があるかを特定する
\end{enumerate}

これらの課題をRmarkdownファイルにまとめ，期日までに提出すること。時間内に終わらなかった場合は，次回授業までに完成させておくこと。分析結果の適切な解釈と可視化を重視し，統計的に有意な結果だけでなく，効果量や信頼区間についても言及するよう心がけること。不明点があれば，\textcite{kosugi2019}や\textcite{Hashimoto201607}を参照すること。